\chapter{Flatness and Descendability}

Let $R$ be a commutative ring and $M$ an $R$-module. Studying the extent to which flat modules are not projective has presented a fruitful 
line of inquiry in homological algebra. 
A truly remarkable characterization of flat modules was given by Raynaud and Gruson in their seminal 1971 paper \cite{Raynaud1971}.
\begin{theorem}\label{thm:rg1}
A module $M$ over a ring $R$ is projective if and only if it is flat, Mittag-Leffler and a direct sum of countably generated modules. 
\cite[Theorem 2.2.1]{Raynaud1971} (see also \cite[Tag 059Z]{sp})
\end{theorem}
The conditions of being Mittag-Leffler and countably generated are related to the well-known fact that the functor:
\[\lim_{I}: \fun(I, \mod{R}) \to \mod{R}\]
has derived functors $R^n\lim_{I}$ whose vanishing is related to the size of the indexing category $I$.
In particular, if $I$ is countable, then $R^k\lim_I = 0$ for all $k \geq 2$, and more generally if $|I|< \aleph_n$
then $R^k\lim_I = 0$ for all $k \geq n+1$. After writing $M$ as a filtered colimit of finite free $R$-modules,
the countably generated assumption allows us to regard $M$ as a filtered colimit indexed by a countable category, 
and the Mittag-Leffler condition ensures a vanishing of $R^1\lim_I$. Theorem~\ref{thm:rg1}
led to a proof that projective modules $M$ satisfy descent with respect to faithfully flat extensions $R \to S$, i.e. if
$M \otimes_R S$ is a projective $S$-module, then $M$ is a projective $R$-module (see \cite[Tag 05A4]{sp}).\\
\indent Theorem~\ref{thm:rg1} has been generalized in various directions, one of which is the relationship between the base ring $R$ and the projective
dimension of a flat module $M$. The following two results are quite remarkable:
\begin{theorem}\label{thm:gj0}
Let $R$ be a commutative ring and $M$ a flat $R$-module.
\begin{enumerate}
\item If $|R| \leq \aleph_n$, then $M$ has projective dimension at most $n+1$. \cite[Theorem 7.10]{gj}
\item If $R$ is Noetherian with $\text{dim}(R) = d$, then $M$ has projective dimension at most $d$. \cite[Corollary 7.2]{gj}
\end{enumerate}
\end{theorem}
These results suggest that there should exist flat modules of infinite projective dimension over rings of large cardinality and infinite Krull dimension.
This suggestion is especially convincing when one remarks that there exists rings $R$, called \textit{absolutely flat rings}, which are characterized by all their modules being flat.
Boolean rings $R$ for example, where for any $r\in R$, $r^2=r$, are absolutely flat, and Osofsky in \cite{osofsky} showed that there exist boolean rings of arbitrarily large cardinality
that admit modules of infinite projective dimension.\\
\indent The techniques used to prove Theorem \ref{thm:gj0} are quite unique, and we wish to briefly summarize them here (see also \cite{jiangdescendable} for an $\infty$-categorical perspective). First 
we recall that in a category $\scr{C}$, an object $X$ is said to be $\aleph_n$-Artinian if any filtered descending chain of subobjects of $X$ 
has a cofinal subchain of length at most $\aleph_n$. The key technical result used to prove Theorem \ref{thm:gj0} is the following:
\begin{theorem}\label{thm:gj1}
Let $\scr{C}$ be an abelian category, and $T:\scr{C} \to \mathrm{Ab}$ an exact functor to the category of abelian groups $\mathrm{Ab}$. If every object of $\scr{C}$ is $\aleph_n$-Artinian,
then the injective dimension of $T$ as an object of the functor category $\mathrm{Lex}(\scr{C}, \mathrm{Ab})$ of left exact functors is at most $n+1$. 
\cite[Theorem 6.5]{gj}
\end{theorem}
We also recall that an exact sequence 
\[\scr{E}: 0 \to A \to B \to C \to 0\]
of $R$-modules is said to be \textit{pure} if for any $R$-module $M$, the induced sequence
\[0 \to A \otimes_R N \to B \otimes_R N \to C \otimes_R N \to 0\]
is exact for any $R$-module $N$. One can observe that $\scr{E}$ is pure if and only if $C$ is a flat $R$-module, if and only if $A \to B$ is \textit{universally injective}.
 A module $K$ is said to be \textit{pure injective} if for any pure exact sequence $\scr{E}$, the induced sequence
\[0 \to \Hom_R(C, K) \to \Hom_R(B, K) \to \Hom_R(A, K) \to 0\]
is exact. Let $\mod{R^{\mathrm{fp}}}$ denote the category of finitely presented $R$-modules. Let $\overline{\overline{M}}$ denote the functor:
\[\overline{\overline{M}}: \mod{R^{\mathrm{fp}}} \to \mathrm{Ab}, \quad  \overline{\overline{M}}(F) = F \otimes_R M\]
for an $R$-module $M$.
\begin{theorem}\label{thm:gj2}
Let $R$ be a commutative ring.
\begin{enumerate}
\item [(i)] A functor $F \in \mathrm{Lex}(\mod{R^{\mathrm{fp}}}, \mathrm{Ab})$ is injective if and only if $F \cong \Hom_R(-, K)$ for some pure injective $R$-module $K$. 
\cite[Proposition 1.2]{gj}
\item [(ii)] If the injective dimension of $\overline{\overline{M}}$ in $\fun(\mod{R^{\mathrm{fp}}}, \mathrm{Ab})$ is at most $n$ for any flat module $M$,
then the projective dimension of any flat $R$-module is at most $n$. \cite[Proposition 3.8]{gj}
\end{enumerate}
\end{theorem}
Finally, to prove Theorem \ref{thm:gj0} (i), one shows that any finitely presented module over a ring $R$ of cardinality at most $\aleph_n$
is $\aleph_n$-Noetherian. We have an equivalence:
\[\Phi: \mathrm{Lex}(\mod{R^{\mathrm{fp}}}, \mathrm{Ab}) \cong \mathrm{Lex}(C(R)^{\text{op}}, \mathrm{Ab}), \; \; F \mapsto \Hom_{\mathrm{Lex}(\mod{R^{\mathrm{fp}}}, \mathrm{Ab})}(-,F)\]
where $C(R)^{\text{op}}$ is the full subcategory of coherent functors in $\mathrm{Lex}(\mod{R^{\mathrm{fp}}}, \mathrm{Ab})$,
and a functor $F$ is coherent if it fits into an exact sequence
\[\Hom_R(-, M) \to \Hom_R(-, N) \to F \to 0\]
for some finitely presented $R$-modules $M$ and $N$. Moreover, $\Phi(\overline{\overline{M}})$ becomes an exact functor (see \cite[Lemma 5.3]{gj}), and the injective dimension
of $\Phi(\overline{\overline{M}})$ and $\overline{\overline{M}}$ in their respective abelian categories is the same. 
Any object in $C(R)$ is $\aleph_n$-Noetherian (see \cite[Theorem 7.9]{gj}), so therefore
every object of $C(R)^{\text{op}}$ is $\aleph_n$-Artinian (see \cite[Theorem 5.6]{gj}), so we can conclude by Theorem~\ref{thm:gj1} and Theorem~\ref{thm:gj2}.\\
\indent In this chapter, we study the relationship between flatness and projective dimension from a different perspective.
In particular, we are interested in the notion of \textit{descendable} map of rings $f: A \to B$, which can be thought of as a strong version of $f$ satisfying descent. 
To be more precise, $f$ is descendable if the map $\{A\}_n \to \{ \text{Tot}_n B^{\bullet}\}_n$ is a pro-isomorphism of $A$-modules where $B^{\bullet}$ 
is the derived Cech nerve and on the other hand, for descent, one usually requires the weaker condition that 
$A \to R\lim B^{\bullet}$ is an isomorphism. 
For a more detailed study of descendability we refer the reader to \cite{DescendMathew}. We briefly mention two results that illustrate the good nature of descendability.
\begin{lemma}
Let $A \to B$ be a descendable map of rings.
\begin{enumerate}
\item[(i)] Let $B \to C$ be a map of rings. If $A \to C$ is descendable, then so is $A \to B$.
\item[(ii)] Let $A \to D$ be a map of rings. Then $D \to B \otimes_A D$ is descendable.
\end{enumerate}
\end{lemma}
\begin{proof}
$A \to B$ is descendable when there exists an $n \in \bb{N}$ such that the map $A \to \lim_{\Delta_{\le n}} B^{\bullet}$ admits a section in the derived category of $A$-modules.
For (i), since $A \to C$ is descendable, there exists an $n \in \bb{N}$ such that $A \to \lim_{\Delta_{\le n}} C^{\bullet}$ admits a section, and we need only pre-compose this section with
the map $\lim_{\Delta_{\le n}} B^{\bullet} \to \lim_{\Delta_{\le n}} C^{\bullet}$ to obtain a section for $A \to \lim_{\Delta_{\le n}} B^{\bullet}$. 
(ii) follows a similar argument by base-changing the section of $A \to \lim_{\Delta_{\le n}} B^{\bullet}$ along $A \to D$.
\end{proof}
It is a celebrated theorem that the $2$-category of quasi-coherent modules have descent with respect to faithfully flat ring maps (see \cite[Tag0306]{sp}).
For descendable maps, we have a much more general result:
\begin{theorem}\label{thm:luriedescend}
Let $A \to B$ be a descendable map of ($\bb{E}_{\infty}$)-rings. Let $\scr{C}$ be a stable $A$-linear $\infty$-category. Then the functor:
\[\scr{C} \to \mathrm{LMod}_{B}(\scr{C}), \quad X \mapsto B \otimes_A X\]
is comonadic, i.e. we have an equivalence of $A$-linear $\infty$-categories:
\[\scr{C} \cong \lim_{B^{\bullet}}(\scr{C})\]
where $B^{\bullet}$ is the derived Cech nerve of $A \to B$. \cite[D.3.4.1]{LurieSAG}
\end{theorem}
In this chapter, we are interested in the question of whether faithfully flat rings maps are \text{descendable}. 
In concrete terms, if $f: A \to B$ is a faithfully flat map of (discrete, commutative) rings, we have an exact sequence $A \to B \to \coker(f)$ 
corresponding to an element $\eta \in \Ext_A^{1}(\coker(f), A)$. The map $f$ will be descendable precisely when there exists an 
integer $n$ such that $\eta^{\otimes_A n}=0$ as an element of $\Ext^n_A(\coker(f)^{\otimes_A n}, A)$.\\
\indent By Theorem~\ref{thm:gj0}, we know that if $A$ is a ring of cardinality at most $\aleph_n$, then any flat $A$-module has projective dimension at most $n+1$,
so therefore $f$ is descendable, of exponent at most $n+1$, since $\coker(f)^{\otimes_A k}$ is a flat $A$-module for any $k > 0$. 
Similarly, if $A$ is a Noetherian ring of Krull dimension $d$, then any flat $A$-module has projective dimension at most $d$,
so $f$ is descendable of exponent at most $d$.\\
\indent Therefore, if faithfully flat rings maps were to fail to be descendable, then the base ring $A$ would have to satisfy $|A| \ge \aleph_{\omega}$ and, 
if Noetherian, $\dim(A) =\infty$. We will construct examples of faithfully flat ring maps $A \to A'$ that is not descendable for a large
class of base rings $A$ that contain an \textit{indivisible sequence}, such as various boolean rings or polynomial rings in infinitely many variables over a `large' field. 
Our strategy can be summarized as follows:
\begin{enumerate}
\item[(i)] Show that there exists rings $A$ with arbitrarily large non-vanishing cup-products, 
i.e. $A$-linear module maps $f_i: M_i \to N_i$ such that $\otimes_{i\in [n],A} f_i \neq 0$, 
using \textit{indivisible sequences} (Definition~\ref{definition:indivisible}).
\item[(ii)] Devise a general procedure to construct non-descendable faithfully flat ring maps from non-vanishing cup-products of universally
injective module maps.
\end{enumerate}
One may wonder whether, granted (ii), one can construct non-descendable modules in a more straightforward manner given any flat module
$M$ of infinite projective dimension. In particular, the author initially speculated that it would be possible to take a flat module $M$ of infinite projective dimension,
 consider a projective resultion $P_{\bullet} \to M$ with boundaries $M_i = \ker(P_i \to P_{i-1})$ (and define $M_0=M$) forming extension classes $\eta_i \in \Ext^1_A(M_i, M_{i+1})$
whose composition $\eta_0 \circ \eta_1 \circ \eta_2 \circ... \circ \eta_n$ would be non-zero for all $n\ge 1$, and then try to construct non-vanishing cup-products from these extension classes.
We believe the following result illustrates why this shouldn't be possible,
since it could be the case that $M_i \otimes_A M_{i+1}=0$ for at least some $i\ge 1$, and therefore 
that descendability is more subtle than just the projective dimension of flat modules.
\begin{lemma}
There exists a ring $A$ and a flat $A$-module $M$ such that $M \otimes_A M = 0$ but $M \neq 0$.
\end{lemma}
\begin{proof}
Let $A=k[x_1, x_2, \ldots]/(x_1^2, x_2^2, \ldots)$ where $k$ is a field. Let $M=\colim_{i \in I} A$ where $I=\bb{N}$ and the 
maps in the colimit are given by multiplication by $x_i$.
Then $M$ is flat since it is a filtered colimit of free modules. However, $M \otimes_A M = 0$ since $M \otimes_A M = \colim_{(i,j) \in I \times I} A$ 
where the maps $(i,j) \to (i+1,j)$ is multiplication by $x_{i+1}$ and the map $(i,j) \to (i,j+1)$ is multiplication by 
$x_{j+1}$, so any element in $A$ in position $(i,j)$ is sent to $0$ in the factor of $A$ in position $(\max(i,j)+1, \max(i,j)+1)$.
\end{proof}

\subsubsection*{Notation and conventions} Let $A$ be a commutative ring. The derived category $D(A)$ of $A$-modules, regarded as a stable $\infty$-category, admits a symmetric monoidal structure given by $\otimesl_A$. Therefore, given $A$-modules $M,N,M',N'$, there is an $A$-bilinear pairing:
\[\text{RHom}_A(M,N) \times \text{RHom}_A(M',N') \to \text{RHom}_A(M \otimesl_A M', N \otimesl_A N').\]
Passing to the homotopy category, we obtain an $A$-bilinear pairing of $A$-modules:
\[\Ext^0_A(M,N) \times \Ext^0_A(M',N') \to \Ext^0_A(M \otimesl_A M', N \otimesl_A N'),\]
and for a pair of $f \in \Ext^0_A(M,N), g \in \Ext^0_A(M',N')$, we will set 
$f \otimesl_A g \in \Ext^0_A(M \otimesl_A M', N \otimesl_A N')$ to be the image of $(f,g)$ under the aforementioned pairing. 
Moreover, if $N=N'=A$, we will compose with the multiplication isomorphism $A \otimesl_A A \to A$ so that the target is 
$\Ext^0_A(M \otimesl_A M', A)$, and note that flatness conditions on our modules remove the necessity for derived tensor products.
\begin{enumerate}
\item[(i)] If $S$ is a set and $R$ a ring, define $\Hom^{\mathrm{fin}}_{\mathrm{Set}}(S, R) \subset \Hom_{\mathrm{Set}}(S,R)$ to be the set-theoretic maps with finite support i.e. $f \in \Hom^{\mathrm{fin}}_{\mathrm{Set}}(S, R)$ if there exists a subset $K \subset S$ with $|K| < \infty$ such that $f(S\setminus K) = 0$. For each $t \in S$, set $\delta_{s=t} \in \Hom^{\mathrm{fin}}_{\mathrm{Set}}(S, R)$ to be the function that is $1$ if $s=t$ and $0$ else.
\item[(ii)]Further, if $S_i$ are sets indexed by the set $\{1,2,...,n\}$, then we will identify the vector $1_{s_1,s_2,...,s_n}$ in the nested direct sum $\oplus_{S_1} \oplus_{S_2} \oplus ... \oplus_{S_n} R$ as follows: Every vector in $v \in \oplus_{S_1} \oplus_{S_2} \oplus ... \oplus_{S_n} R$ can be identified as a function $f_{v} \in \Hom^{\mathrm{fin}}_{\mathrm{Set}}(\prod_i S_i, R)$, and under this identification, $1_{s_1,...,s_n}$ corresponds to $\delta_{t=s}$ where $s=(s_1,s_2,...,s_n) \in \prod_i S_i$.
\item[(iii)]For any set $S$, we let $p^n_i: S^{\times n} \to S^{\times n-1}$ be defined as follows: For any $s = (s_1,s_2,..,s_n) \in S^{\times n}$,
\[p^n_i(s) = (s_1,s_2,...,s_{i-1},s_{i+1},...,s_n).\]
For a point $s = (s_1,...,s_{n-1}) \in S^{\times n-1}$, we consider the degeneracy map $t_i^{s}: S \to S^{\times n}$ which takes a point $s' \in S$ to the point:
\[(s_1, ..., s_{i-1}, s', s_{i},...,s_{n-1}) \in S^{\times n}.\]
\item[(iv)] For a ring $R$ and a set $S$, let $R[S]=\Sym_R(\oplus_{s \in S} R) = R[X_s]_{s \in S}$ and $R\{S\} = R[X_s]_{s \in S}/I$ where $I$ is the ideal generated by $X_s^2 - X_s \in R[S]$. \\
\indent Let $B(S) = \{K \subset S| \; |K| < \infty\}$, which can be endowed with the structure of a group under $\cup$, with unit $\emptyset$, which moreover satisfies $x \cup x = x$ for any $x \in B(S)$. Then $R\{S\} \simeq R[B(S)]$, the group ring on $B(S)$, and so we have an $R$-linear decomposition $R\{S\} = \oplus_n R\{S\}_{=n}$, where $R\{S\}_{=n} = \oplus_{\underset{|K|=n}{K\subset S}} R$. Let $\iota_S: S\to B(S)$ be the set-theoretic inclusion viewing elements of $S$ as one-element subsets of $S$, and let $R(\iota_S): \oplus_S R \to R\{S\}$ be the corresponding identification of $\oplus_S R$ as the degree $1$ elements of $R\{S\}$. \\
\indent We remark that $B(S \amalg T) = B(S) \times B(T)$ and therefore $R\{S \amalg T\} \simeq R\{S\} \otimes_R R\{T\}$.
\item[(v)] For an (injective) $A$-linear map $f: M \to N$ of $A$-modules, we may form an exact sequence 
\[0 \to M \to N \to \coker{f} \to 0.\]
This corresponds to a class in $\Ext^1_A(\coker{f}, M)$, which we denote by $\cl{f}$.\\
\indent For a ring map $f: A \to B$ that is faithfully flat, we say that $f$ has \textit{descendability exponent $n \in \bb{N}$} 
when $n$ is the minimal number such that $\text{cl}(f)^{\otimes_A n} \neq 0$ but $\text{cl}(f)^{\otimes_A n+1} = 0$. 
When $n$ doesn't exist, we say that $f$ is \textit{not descendable}, and has descendability exponent $\infty$.\footnote{In the literature one will typically define the exponent of descendability to be the minimal $n$ such that $\text{cl}(f)^{\otimes_A n} = 0$, 
however, due to the relationship between descendability and projective dimension of flat modules, we find this definition more convenient for our purposes. 
So our definition of exponent of descendability is always one less than the usual definition.}
\end{enumerate}
For any natural number $n \in \bb{N}$ greater than $1$, let $[n] = \{1,2,...,n\}$ and set $[\infty] = \bb{N}_{> 0}$.

\section{Indivisible Sequences and Descendability}

\begin{definition}\label{definition:indivisible}
Let $R$ be a commutative ring. For a set $S$, and a set-theoretic map
\[\Psi: S \to R\]
we associate the $R$-linear map
\[\Psi_{R,S}: \oplus_{S} R \to R \oplus \bigoplus_{S} R\]
defined by $\Psi_{R,S}(1_s)_{0}= 1$ and $\Psi_{R,S}(1_{s})_{s'} = \delta_{s'=s}$, where we regard the factor not indexed by $S$ on the right-hand side as being in degree $0$.\\
\indent A \textit{$1$-indivisible sequence} is the data $\{S, \Psi: S \to R\}$ where $S$ is a set and $\Psi:S \to R$ a set-theoretic map such that the associated map $\Psi_{R,S}$ is an 
\textit{$R$-universally injective map}, in the sense that $\Psi_{R,S} \otimes_{R,f}R'$ is injective for any ring map $R \over{f}{\to} R'$. Equivalently, $\Psi_{R,S}$ is injective and 
$\coker{\Psi_{R,S}}$ is $R$-flat (as $\text{Tor}^1_R(\coker{\Psi_{R,S}}, R')=0$ for any ring map $R \to R'$).\\
\indent For any $n \in \bb{N} \cup \{\infty\}$, an \textit{$n$-indivisible sequence} is the data $\{\{S_i, \Psi_i:S_i \to R\}\}_{i \in [n]}$ where each $\{S_i, \Psi_i: S_i\to R\}$ 
is a $1$-indivisible sequence, and the elements $\{\Psi_i(s_i)\}_{i \in [n]}$ of $R$ do not generate the unit ideal for any choice of 
elements $s_i \in S_i$. An $\infty$-indivisible sequence of $R$ will simply be called an \textit{indivisible sequence} of $R$.
\end{definition}
While it may not be immediately obvious, there are many examples of rings $R$ with $n$-indivisible sequences. First, there is a general way to produce $n$-indivisible sequences from $1$-indivisible sequences.
\begin{proposition}\label{proposition:indivequiv}
Let $k$ be a commutative ring and $f: R \to R'$ a map of $k$-algebras, and for any $i \in [n]$, let $t_i: R \to R^{\otimes_k n}$ be the inclusion by $1$ into the $i^{\text{th}}$ factor.
\begin{enumerate}
\item[(i)] $\Psi_{R,S} \otimes_{R,f} R' = (f \circ \Psi)_{R',S}$. In particular, if $\{S, \Psi: S \to R\}$ is $1$-indivisible, then so is $\{S, f \circ \Psi: S \to R'\}$.
If $f$ is faithfully flat, then $\{\{S_i, \Psi: S \to R\}\}$ is $n$-indivisible if and only if $\{\{S_i, f \circ \Psi: S \to R'\}\}$ is $n$-indivisible.
\item[(ii)] For each $n \in \bb{N} \cup \{\infty\}$ and $i \in [n]$, consider $1$-indivisible sequences $\{S_i, \Psi_i:S_i \to R\}$ such that for every $s_i \in S_i$, the composed map 
\[k \to R \to R/\Psi_i(s_i)\]
is faithfully flat. Then $\{\{S_i, t_i \circ \Psi_i: S_i \to R^{\otimes_k n}\}\}_{i \in [n]}$ is an $n$-indivisible sequence.
\end{enumerate}
\end{proposition}
\begin{proof}
(i) is standard. For (ii), by (i) the sequences $\{S_i, t_i \circ \Psi_i:S_i \to R^{\otimes_k n}\}$ are $1$-indivisible for each $i \in [n]$. 
Using the K\"unneth spectral sequence applied to the complexes $R \over{\Psi(s_i)}{\to} R$, we can show that:
\[R^{\otimes_{k} n}/(t_1 \circ \Psi_1({s_1}),...,t_n \circ \Psi_n(s_n)) = \bigotimes_{i\in [n], k} R/\Psi_i(s_i).\]
for arbitrary elements $s_i \in S_i$. Under the faithfully flat assumption of the rings on the right, the latter ring is non-zero, hence (ii).
\end{proof}
\begin{example}\label{example:indiv}
Here we will list a few examples of $n$-indivisible sequences.
\begin{enumerate}
\item[(i)] Let $k$ be a commutative ring. Then $\{k, \Psi:k \to k[x]\}$, where $\Psi(a) = x-a$ for each $a \in k$, 
is a $1$-indivisible sequence. Using Proposition~\ref{proposition:indivequiv} we find that 
$\{\{k, \Psi_i: k \to k[x_1,...,x_n]\}\}_{i\in [n]}$, where $\Psi_i(a) = x_i - a$ for every $a \in k$, forms an $n$-indivisible sequence, 
for any $n \in \bb{N} \cup \{\infty\}$. 
\item[(ii)] Let $S$ be a set and $R$ be a ring with a subset $\{e_s\}_{s \in S}$ of orthogonal idempotents. Then the sequence 
$\{S, \Psi:S \to R\}$ defined by $\Psi(s) = 1-e_s$ is $1$-indivisible. Using Proposition~\ref{proposition:indivequiv}, we find that 
$\{\{S, t_i \circ \Psi:S \to R^{\otimes_{\bb{F}_p} n}\}\}_{i \in [n]}$, forms an $n$-indivisible sequence for any 
$n \in \bb{N}\cup \{\infty\}$. We remark that $R^{\otimes_{\bb{F}_p} n}$ is a $p$-boolean ring (see \ref{definition:pboolean}).
\end{enumerate}
Let us now briefly discuss why each of the examples are $1$-indivisible. For (i), we note that
\[\coker{\Psi_{k[x],k}} = \sum_{a \in k} \frac{1}{x-a}k[x] \subset \text{Frac}(k[x]).\]
It's standard to show that $\coker{\Psi_{k[x],k}}$ is therefore flat. Since $k[x]$ is a domain, $\Psi_{k[x],k}$ is clearly injective, 
so $\Psi_{k[x],k}$ is universally injective.\\
\indent For (ii), by (i) Proposition~\ref{proposition:indivequiv}, it suffices to show that $\Psi_{R,S}$ is injective. 
In particular, suppose we have an equation:
\[\Psi_{R,S}(\sum_{i \in [n]} r_i 1_{s_i}) = (\sum_{i \in [n]} r_i) \oplus \bigoplus_{i \in [n]} \Psi(s_i)r_i = 0\]
for a subset of distinct elements $\{s_1,...,s_n\} \subset S$ and $r_i \in R$. Therefore, $r_i \in e_{s_i}R, \; \; \forall i \in [n]$, 
and 
\[\sum_{i \in [n]} r_i = 0,\]
which, after multiplying by $e_{s_i}$, gives that $r_ie_{s_i}=0,$ and therefore $r_i \in (1-e_{s_i})R \Rightarrow r_i=0 \;\; \forall i \in [n]$, as desired. 
\end{example}
The key property of $n$-indivisible sequences is that they give rise to non-vanishing cup-products in $\Ext^n$ groups.
\begin{theorem}\label{theorem:indivisiblecup}
Let $n \in \bb{N}$, $R$ be a ring and $\{\{S_i, \Psi_i: S_i \to R\}\}_{i \in [n]}$ an $n$-indivisible sequence in $R$. Then the class:
\[\bigotimes_{i \in [n],R} (R(\iota_{S_i}) \circ \cl{\Psi_{R,S_i}}) \in \Ext^{n}_{R}(\otimes_{i \in [n],R} \coker{\Psi_{R,S_i}}, R\{\underset{i \in [n]}{\amalg} S_i\})\]
is non-zero as soon as $|S_i| \ge \aleph_{i-1}, \; \forall i \in [n]$.\footnote{Note that:
\[\bigotimes_{i\in [n],R}^n (R(\iota_{S_i}) \circ \cl{\Psi_{R,S_i}}) \neq 0 \Longleftrightarrow \bigotimes_{i\in [n],R} (\cl{\Psi_{R,S_i}}) \neq 0\]
since for each $i \in [n]$, $R({\iota}_{S_i}): \oplus_{S_i} R \to R\{S_i\}$ admits an $R$-linear left-inverse.}
\end{theorem}
The following lemma is helpful for showing that certain $\Ext^n$ classes are non-zero.
\begin{lemma}\label{lemma:extcrit}
Let $A$ be a ring, $P_1, P_2, P_1', P_2', ..., P_r'$ be projective $A$-modules, together with two flat $A$-modules $M$ and $M'$ that form two exact sequences:
\[0 \to P_1 \stackrel{d_1}{\to} P_2 \to M \to 0,\]
\[0 \to P_1' \stackrel{d_2}{\to} ... \to P_r' \to M' \to 0.\]
We may view these as extensions $\eta_1 \in \Ext^1_A(M, P_1)$ and $\eta_2 \in \Ext^{r-1}_A(M',P_1')$, 
and as projective resolutions $P^{\bullet}_1, P^{\bullet}_2$. Then the projective resolution:
\[P^{\bullet}_1 \otimes_A P^{\bullet}_2 \to M \otimes_A M',\]
gives rise to an extension $\Ext^{r}_A(M \otimes_A M', P_1 \otimes_A P_1')$ which is equal to $\eta_1 \otimes_A \eta_2$. 
Moreover, showing this extension is non-zero is equivalent to verifying that 
\[P_1 \otimes P_1' \stackrel{d_1 \otimes 1_{P'_1} \oplus 1_{P_1} \otimes d_2}{\to} P_2 \otimes_A P_1' \oplus P_1 \otimes_A P_2\]
doesn't admit an $A$-linear left inverse.
\end{lemma}
\begin{proof}
For the first claim, the resolutions $P_1^{\bullet}, P_2^{\bullet}$ are cofibrant replacements for $M$ and $M'$, and the 
maps $P_1^{\bullet}[-(r-1)] \to P_1$ and $P_2^{\bullet}[-1] \to P_1'$ correspond to chain maps that are identities on degree 0. 
The tensor product of this map is then just the identity on $P_1 \otimes_A P_1'$ in $P_1^{\bullet} \otimes_A P_2^{\bullet}[-r]$, 
which corresponds to the extension class of $M\otimes_A M'$ by $P_1 \otimes_A P_1'$ given by $P_1^{\bullet} \otimes_A P_2^{\bullet}$. 
The final claim then follows readily, as $\eta_1 \otimes_A \eta_2 = 0$ if and only $\mathrm{id}_{P_1 \otimes_A P_1'}$ is in the image 
of the map 
\[\Hom_A(P_2 \otimes_A P_1' \oplus P_1 \otimes_A P_2, P_1 \otimes_A P_1') \to \Hom_A(P_1 \otimes_A P_1', P_1 \otimes_A P_1')\]
induced by $d_1 \otimes 1_{P'_1} \oplus 1_{P_1} \otimes d_2$.
\end{proof}
\begin{proof}[Proof of Theorem~\ref{theorem:indivisiblecup}]
Using Lemma~\ref{lemma:extcrit}, we may reduce to showing there does not exist an $R$-linear extension $r$ fitting into the 
commutative diagram:
\[\begin{tikzcd}R\{\underset{i \in [n]}{\amalg} S_i\} &[5em]\\[5ex] \underset{{\underset{i \in [n]}{\prod} S_i}}{\bigoplus} R \arrow[r, "\bigoplus_i \Psi_{R,S_i}^n"] \arrow[u,"\underset{i \in [n],R}{\bigotimes} R(\iota_{S_i})"]& \bigoplus_{i} \underset{\underset{j \in [n] \setminus \{i\}}{\prod} S_j}{\bigoplus} R \oplus \underset{\underset{j \in [n]}{\prod}S_j}{\bigoplus}  R \arrow[ul, dashed, swap, "r"]\end{tikzcd}\]
Here, $\Psi^n_{R, S_i}$ is defined as the unique $R$-linear map satisfying:
\[\Psi^n_{R, S_i}(1_s) = 1_{p^n_i(s)} \oplus \Psi_i(s_i) 1_{s}, \; \; \forall s \in \prod_{j \in [n]} S_j.\]
Suppose such an extension $r$ existed. For every $ i \in [n]$, $s' \in \prod_{j \in [n]\setminus \{i\}} S_j$ and 
$s \in \prod_{j \in [n]}S_j$, let
\[r(1_{i,s'}) = f_i(s'), r(1_{i,s}) = g_i(s).\]
Here, $f_i(s'), g_i(s)$ are elements $R\{\amalg_{i \in [n]} S_i\}$, which satisfy the following equation in 
$R\{\amalg_{i \in [n]} S_i\}$: 
\[\sum_i f_i(p^n_i(t)) + \sum_i \Psi_{i}(\text{pr}_i(t)) g_i(t) = \prod_{i \in [n]} R(\iota_{S_i})(1_{\text{pr}_i(t)}), \; \; (*)\]
for every $t \in \prod_{j \in [n]} S_j$.\\
\indent Let 
\[d_n: R\{\amalg_{i \in [n]} S_i\} \to \underset{\underset{i \in [n]}{\prod}S_i}{\bigoplus}  R\]
be the $R$-linear left-inverse to ${\otimes}_{i \in [n],R} R(\iota_{S_i})$. Then for any $t' \in  \prod_{j \in [n]\setminus \{i\}} S_j$, $d_n(f_i(t'))$ 
is a vector in $\bigoplus_{\underset{i \in [n]}{\prod}S_i}  R$, which we view as an element of $\Hom^{\mathrm{fin}}_{\mathrm{Set}}({\prod}_{i \in [n]}S_i , R)$. 
By Proposition~\ref{proposition:setbounds}, there exists an element $s \in \prod_{i \in [n]} S_i$ such that $d_n(f_i(p^n_i(s)))(s)=0, \; \forall i \in [n]$. 
Therefore, applying $d_n$ to $(*)$, we see that:
\[\sum_i \Psi_i(\text{pr}_i(s)) d_n(g_i(s))(s) = d_n(\prod_{i \in [n]} R(\iota_{S_i})(1_{\text{pr}_i(s)}))(s) = 1.\]
But this implies that $\{\Psi_i(\text{pr}_i(s))\}_{i \in [n]}$ generate the unit ideal in $R$, which is a contradiction.
\end{proof}
The following set-theoretic result was used above.
\begin{proposition}\label{proposition:setbounds}
Let $R$ be a ring and $S$ a set. For each $ i\in \{1,...,n\}$, consider functions:
\[v_i: \prod_{j \in [n] \setminus \{i\}} S_j \to \Hom^{\mathrm{fin}}_{\mathrm{Set}}(\prod_{i \in [n]} S_i,R).\]
If $|S_i| \ge \aleph_{i-1}$, then there exists an element $s \in \prod_{i \in [n]} S_i$ such that $v_i(p^n_i(s))(s) = 0$ 
for each $i \in \{1,2,...,n\}$.
\end{proposition}
\begin{proof}
Define $v'_i = v_i(p^n_i(s))(s): \prod_{i \in [n]} S_i \to R$. The functions $v'_i$ have the property that for any 
$s' \in \prod_{j \in  [n] \setminus \{i\}} S_i$, 
\[v'_i \circ t_i^{s'}=v_i(s') \circ t_i^{s'} \in \Hom^{\mathrm{fin}}_{\mathrm{Set}}(S,R).\]
For each $s' \in \prod_{i \in [n-1]} S_i$, define a subset $V_{s'} \subset S_n$ as follows:
\[V_{s'} = \{s \in S_n | \; v'_n(t_n^{s'}(s)) \neq 0 \}.\]
By supposition, $V_{s'}$ is a finite subset of $S_n$. Notice that:
\[|\cup_{s' \in \prod_{i \in [n-1]} S_i} V_{s'}| \le \aleph_0 \aleph_{n-2} < \aleph_{n-1}.\]
Therefore, there is an $s_{n} \in S_n$ such that:
\[v'_n(t_n^{s'}(s_n)) = 0, \; \; \forall s' \in \prod_{i \in [n-1]} S_i.\]
We may repeat this argument for the restricted functions $v'_i \restriction \prod_{j \in [n] \setminus \{i\}} S_j \times s_n$, 
and we eventually end up with elements $s_i \in S_i$ defining an element $s \in \prod_{i \in [n]}S_i$ such that 
$v'_i(s) = v_i(p^n_i(s))(s)=0, \; \; \forall i \in [n]$, as desired.
\end{proof}
\section{Constructing non-descendable ring maps}
Theorem~\ref{theorem:indivisiblecup} allows us to construct non-vanishing cup-products of maps between modules from indivisible sequences.
Given a ring $R$ and a module $M$, there are many possible functorial ways to construct a corresponding algebra, such as $\Sym_R(M)$. 
For any such construction, one would need to be able to carry over the non-vanishing results that we've already proved. The next sequence of results
give a general conditions for which this is possible.
\begin{proposition}\label{proposition:modalg}
Let $n \in \bb{N} \cup \{\infty\}$ and $R$ a commutative ring. Suppose for each $i \in [n]$, we have universally injective $R$-linear maps
 $\Psi_i: M_i \to N_i$, $R$-linear maps $d_i: M_i \to T_i$ with $T_i$ having the structure of a flat $R$-algebra with unit $R \to T_i$ 
 admitting an $R$-linear left-inverse, 
 and $R$-algebra maps 
 $T_i \over{g_i}{\leftarrow} A_i \over{f_i}{\to} B_i$ fitting into a commutative diagram:
\[\begin{tikzcd}
T_i&\\
A_i \arrow[u, "g_i", swap] \arrow[r, "f_i"] & B_i \\
M_i \arrow[u] \arrow[uu, bend left, "d_i"] \arrow[r, "\Psi_{i}"] & N_i \arrow[u]
\end{tikzcd}\]
with each $f_i$ faithfully flat. Suppose that for every finite subset $S' \subset [n]$,
\[\bigotimes_{s'\in S'} (d_{s'} \circ \cl{\Psi_{s'}}) \neq 0.\]
Then the faithfully flat ring map:
\[\begin{tikzcd}\bigotimes_{i \in [n], R} A_i \arrow[r, "\otimes_{i \in [n],R} f_i"]& \bigotimes_{i \in [n],R} B_i\end{tikzcd}\]
has descendability exponent $n$.
\end{proposition}
\begin{proof}
Let $h_i: T_i \to U_i$ be the base-change of $f_i$ along $g_i$. By base-change along 
$\otimes_{i \in [n], R} g_i : \otimes_{i \in [n], R} A_i \to \otimes_{i \in [n], R} T_i$, 
it suffices to show that $\otimes_{i \in [n], R} h_i$ has descendability exponent $n$. But this follows from
Proposition~\ref{proposition:cup} (ii).
\end{proof}
The following proposition was used above.
\begin{proposition}\label{proposition:cup}
Let $A \to B$ be a flat map of rings.
\begin{enumerate}
\item[(i)] For any natural number $n$, let $M_i$ be $B$-modules in $D(B)$ for $i \in [n]$, $\mu_{B/A}^{k}: B^{\otimes_A k+1} \to B$ 
be the multiplication map realizing $B$ as an $A$-algebra, and $\R_{B|A}:D(B) \to D(A)$ be the restriction of scalars functor. 
Consider elements $\eta_i \in \Ext^1_B(M_i, B)$, then:
\[\mu^{n-1}_{B/A} \circ \otimesl_{i\in [n],A} \R_{B|A}(\eta_i) \neq 0 \implies \otimesl_{i\in [n], B} \eta_i \neq 0.\]
\item[(ii)] For every $s \in \bb{N}$, suppose we have flat $A$-algebras $B_s$, faithfully flat maps $f_s: B_s \to C_s$ of $A$-algebras 
and universally injective maps $\Psi_s: M_s \to N_s$ of flat $A$-modules fitting into a commutative diagram:
\[\begin{tikzcd}
B_s \arrow[r, "f_s"] & C_s \\
M_s \arrow[u, "h_s"] \arrow[r, "\Psi_s"] & N_s \arrow[u]
\end{tikzcd}\]
Assume the unit $A \to B_s$ admits an $A$-linear left-inverse $k_s: B_s \to A$. Then for any subset $S \subset \bb{N}$, and any finite subset $S' \subset S$:
\[\bigotimes_{s' \in S',A} (h_{s'} \circ \cl{\Psi_{s'} }) \neq 0 \implies \underset{s' \in S', \otimes_{s \in S,A}B_s}{\bigotimes} \cl{ \otimes_{s \in S,A} f_s} \neq 0.\]
\end{enumerate}
\end{proposition}
\begin{proof}
For (i), we remark more generally that for any pair of morphisms $f: N_1 \to K_1, g: N_2 \to K_2$ in $D(B)$, 
the bar construction of \cite[Chapter 4.4.2]{LurieHA} gives a homotopy coherent diagram:
\[\begin{tikzcd}
N_1 \otimesl_B K_1 \arrow[r, "f \otimesl_B g"] & N_2 \otimesl_B K_2 \\
N_1 \otimesl_A K_1 \arrow[u] \arrow[r, "f \otimesl_A g"] & N_2 \otimesl_A K_2 \arrow[u]
\end{tikzcd}\]
Applying this to the $\eta_i$ repeatedly, we end up with a homotopy commutative diagram:
\[\begin{tikzcd}
\otimesl_{i\in [n], B} M_i[-1] \arrow[r, "\otimesl_{i\in [n],B} \eta_i"] & B^{\otimesl_B n} \arrow[r, "\simeq"] & B \\
\otimesl_{i\in [n], A} M_i[-1] \arrow[u, "\varphi"] \arrow[r,swap, "\otimesl_{i\in [n],B} \R_{B|A}(\eta_i)", {yshift=-3pt}] & B^{\otimesl_A n} \arrow[u] \arrow[ru, swap,"\mu^{n-1}_{B/A}"] &
\end{tikzcd}\]
We therefore see that:
\[\otimesl_{i\in [n],B} \eta_i =0 \implies (\otimesl_{i\in [n],B} \eta_i) \circ \varphi =0 \implies \mu^{n-1}_{B/A} \circ \otimesl_{i\in [n],A} \R_{B|A}(\eta_i) = 0.\]
For (ii), let $B_{S} = \otimes_{s \in S} B_{s}$,  $i_s: B_s \to B_S$ be the inclusion in the $s^{\text{th}}$-factor, and $g_s: \coker{f_s} \to \coker{\otimes_{s \in S} f_s}$ the induced map. Note that:
\[\bigotimes_{s' \in S',A} (h_{s'} \circ \cl{\Psi_{s'} }) \neq 0 \implies \bigotimes_{s' \in S',A} \cl{f_{s'}} \neq 0.\]
By (i) it suffices to show that:
\[\mu^{|S'|-1}_{B_{S}/A} \circ \bigotimes_{s' \in S', A} \R_{B_{S}|A}(\cl{\otimes_{s \in S} f_s}) \neq 0.\]
This follows from the following calculation:
\begin{align*}k_{S'}\! \circ \!\mu^{|S'|-1}_{B_{S}/A} \!\circ \!\bigotimes_{s' \in S', A} \R_{B_{S}|A}(\cl{\otimes_{s \in S} f_s}) \! \circ \!\bigotimes_{s' \in S', A} g_{s'}\! &= \! k_{S'} \!\circ \! \mu^{|S'|-1}_{B_{S}/A} \circ \!\bigotimes_{s' \in S', A} (\R_{B_{S}|A}(\cl{\otimes_{s \in S} f_s}) \circ g_{s'})\!\\
& = k_{S'} \circ \mu^{|S'|-1}_{B_{S}/A} \circ \bigotimes_{s' \in S'} (i_{s'} \circ \cl{f_{s'}})\\
&= \left (k_{S'} \circ \mu^{|S'|-1}_{B_{S}/A} \circ \bigotimes_{s' \in S'} i_{s'}\right) \circ \bigotimes_{s' \in S'} \cl{f_{s'}}\\
&= \bigotimes_{s' \in S'} \cl{f_{s'}} \\
&\neq 0,
\end{align*}
where $k_{S'}: \otimes_{s \in S} R \to \otimes_{s \in S \setminus S'}R$ is defined as the tensor product of 
the $A$-linear maps $\text{id}_{B_{s''}}: B_{s''} \to B_{s''}$ and $k_{s'}$ for each $s'' \in S \setminus S'$ and $s' \in S'$. \\
\end{proof}
\begin{theorem}\label{theorem:indivalg}
Let $n \in \bb{N} \cup \{\infty\}$, and $\{\{S_i, \Psi_i:S_i \to R\}\}$ be a $n$-indivisible sequence in $R$. For all $i \in [n]$, suppose we have $R$-algebra maps $R\{S_i\} \over{g_i}{\leftarrow} A_i \over{f_i}{\to} B_i$ fitting into a commutative diagram:
\[\begin{tikzcd}
R\{S_i\}&\\
A_i \arrow[u, "g_i", swap] \arrow[r, "f_i"] & B_i \\
\oplus_{S_i} R \arrow[u] \arrow[uu, bend left, "R(\iota_{S_i})"] \arrow[r, "\Psi_{R,S_i}"] & R \oplus \bigoplus_{S_i} R \arrow[u]
\end{tikzcd}\]
and each $f_i$ are faithfully flat. If for all $i \in [n]$, $|S_i| \ge \aleph_{i-1}$, then the faithfully flat ring map:
\[\begin{tikzcd}\bigotimes_{i \in [n], R} A_i \arrow[r, "\otimes_{i \in [n],R} f_i"]& \bigotimes_{i \in [n],R} B_i\end{tikzcd}\]
has descendability exponent $n$.
\end{theorem}
\begin{proof}
This follows directly from Theorem~\ref{theorem:indivisiblecup} and Proposition~\ref{proposition:modalg}.
\end{proof}
Before stating our main result, we need to recall a few definitions.
\begin{definition}\label{definition:pboolean}
A ring $R$ is said to be \textit{p-boolean} if $R$ is in characteristic $p$ and for every $r \in R, r^p=r$.
\end{definition}
The following is a discrete version of a derived result appearing in \cite[Theorem 4.6]{pbool}.
\begin{theorem}\label{theorem:p-booleanmonadic}
Let $R$ be a $p$-boolean ring. Let $\mathrm{CAlg}_{R}$ (resp. $\mathrm{CAlg}^{\mathrm{perf}}_{R}$, $\mathrm{CAlg}^{\phi=1}_{R}$) denote the category of (commutative) $R$-algebras (resp. perfect $R$-algebras, $p$-boolean $R$-algebras). Each of the forgetful functors:
\[\mathrm{CAlg}^{\phi=1}_{R} \to \mathrm{CAlg}^{\mathrm{perf}}_{R} \to \mathrm{CAlg}_{R} \to \text{Mod}_R\]
are monadic. In particular, the left adjoint of this forgetful functor is the functor $F_R^p: \text{Mod}_R \to \mathrm{CAlg}^{\phi=1}_{R}$ defined on objects by the rule:
\[M \mapsto \Sym_R(M) \mapsto \Sym_R(M)_{\mathrm{perf}} \mapsto \coeq(\phi, \mathrm{id}: \Sym_R(M)_{\mathrm{perf}} \to \Sym_R(M)_{\mathrm{perf}}),\]
where $(-)_{\mathrm{perf}}$ is the colimit perfection functor of commutative algebras that on objects is defined as:
\[A \mapsto A_{\mathrm{perf}} = \colim(A \stackrel{\phi}{\to} A \stackrel{\phi}{\to} A \stackrel{\phi}{\to} ...),\]
and $\phi$ is the Frobenius endomorphism of $A$.
\end{theorem}
\begin{proposition}\label{proposition:symformula}
Let $R$ be a ring and $f: M \to N$ an injective map of flat $R$-modules such that $K=\coker{f}$ is $R$-flat. Then 
\[\Sym_R(f): \Sym_R(M) \to \Sym_R(N)\]
is a faithfully flat map of $R$-algebras. Furthermore, if $R$ is $p$-boolean, then $F_R^p(f)$ is a faithfully flat map of $p$-boolean rings.
\end{proposition}
\begin{proof}
By Lazard's theorem, $K=\colim_{i\in I} K_i$ where $I$ is a filtered indexing category and $K_i$ is a finite free $R$-module. If we define $N_i = N \times_{K} K_i$, then
\[\colim_{i \in I} N_i = \colim_{i \in I} N \times_{K} K_i = N \times_{K} \colim_{i \in I} K_i = N,\]
and we have an exact sequence
\[0 \to M \over{f_i}{\to} N_i \to K_i \to 0,\]
which is split. Therefore:
\[\Sym_R(N_i) = \Sym_R(M) \otimes_R \Sym_R(K_i)\]
and if $R$ is $p$-boolean:
\[F^p_R(N_i) = F_R^p(M) \otimes_R F_R^p(K_i).\]
Note that as $K_i$ are finite free, the unit $R \to \Sym_R(K_i)$ is faithfully flat, and if $R$ is futher assumed to be $p$-boolean, then likewise for $R \to F^p_R(K_i)$.\footnote{If $K_i =R^{\oplus n}$, then $F^p_R(K_i) = R[X_1,...,X_n]/(X_1^p-X_1, ..., X_n^p-X_n)$.} Hence, the map:
\[\Sym_R(M) \over{\Sym_R(f_i)}{\to} \Sym_R(N_i)\]
is faithfully flat, and if $R$ is $p$-boolean, then:
\[F^p_R(M) \over{F^p_R(f_i)}{\to} F^p_R(N_i)\]
is also faithfully flat. The result follows by noting $\Sym_R(f) = \colim_{i\in I}\Sym_R(f_i)$ and if $R$ is $p$-boolean, $F^p_R(f) = \colim_{i \in I} F^p_R(f_i)$.
\end{proof}
Our main result on descendability of faithfully flat ring maps can then be stated as follows.
\begin{theorem}\label{theorem:mainindiv}
Let $n \in \bb{N} \cup \{\infty\}$, $R$ be a commutative ring and $\{\{S_i, \Psi_i: S_i \to R\}\}_{i \in [n]}$ an $n$-indivisible sequence, and $|S_i| \ge \aleph_{i-1}$. Then the ring map 
\[\bigotimes_{i=1, R}^n \Sym_R(\Psi_{R,S_i}): \bigotimes_{i=1, R}^n \Sym_R(\oplus_{S_i} R) \to \bigotimes_{i=1, R}^n \Sym_R(R \oplus \bigoplus_{S_i} R),\]
is faithfully flat with exponent of descendability $n$. In fact, the ring map:
\[\begin{tikzcd}R\{\underset{i=1}{\amalg^n}S_i\}\arrow[r, "\underset{i=1,R}{\bigotimes^n} \Sym_R(\Psi_{R,S_i}) \underset{\underset{i=1,R}{\bigotimes^n} \Sym_R(\oplus_{S_i} R)}{\bigotimes}  R\{\underset{i=1}{\amalg^n}S_i\}"] &[8em] \underset{i=1,R}{\bigotimes^n} \Sym_R(R \oplus \bigoplus_{S_i} R)  \underset{\underset{i=1,R}{\bigotimes^n} \Sym_R(\oplus_{S_i} R)}{\bigotimes}R\{\underset{i=1}{\amalg^n}S_i\}\end{tikzcd},\]
also has descendability exponent $n$ too.\\
\indent If $R$ is $p$-boolean, then the ring map:
\[\bigotimes_{i=1, R}^n F^p_R(\Psi_{R,S_i}): \bigotimes_{i=1, R}^n F^p_R(\oplus_{S_i} R) \to \bigotimes_{i=1, R}^n F^p_R(R \oplus \bigoplus_{S_i} R),\]
is a faithfully flat of $p$-boolean rings of descendablity exponent $n$.
\end{theorem}
\begin{proof}
After noting commutative diagrams:
\[\begin{tikzcd} \oplus_{S_i} R \arrow[r] \arrow[rd]\arrow[rr, bend left, "R(\iota_{S_i})"]&  \Sym_R(\oplus_{S_i} R) \arrow[r]& R\{S_i\}\\
& F^p_R(\oplus_{S_i} R) \arrow[ru]&
\end{tikzcd}\]
the result follows from combining Proposition~\ref{proposition:symformula} and Theorem~\ref{theorem:indivalg}.
\end{proof}
\begin{corollary}\label{corollary:mainindiv}
For any $n \in \bb{N} \cup \{\infty\}$:
\begin{enumerate}
\item[(i)] There exists a faithfully flat ring map $A \to A'$ between $\aleph_{n-1}$-Noetherian rings, with $A$ of Krull-dimension $n$, which has descendability exponent $n$. 
\item[(ii)] There exists a faithfully flat ring map $A \to A'$ between $\text{min}(\beth_{n-1}, 2^{\aleph_{n-1}})$-countable $p$-boolean rings which has descendability exponent $n$.
\end{enumerate}
\end{corollary}
\begin{proof}
For (i), we may use the first example of Example~\ref{example:indiv}. In particular, if we set $R = k[x_1,x_2,...,x_n]$, 
then from Theorem~\ref{theorem:mainindiv} we may conclude that there is a faithfully flat map:
\[k\{k^{\amalg n}\}[x_1,...,x_n] = R\{k^{\amalg n}\} \to R\{k_{+}^{\amalg n}\}= A'\]
which has descendability exponent $n$ when $|k| \ge \aleph_{n-1}$, where the set $k_+ = \{*\}\amalg k$. 
Furthermore, it's clear that $R\{k^{\amalg n}\}$ has Krull-dimension $n$ (its irreducible components are isomorphic to R), 
and that the rings are $\aleph_{n-1}$-Noetherian (since they're $\aleph_n$-countable).\\
\indent For (ii), by using Theorem~\ref{theorem:mainindiv} and Proposition~\ref{proposition:indivequiv}, 
it suffices to find a $1$-indivisible sequence $\Psi: S \to R$ where the image of $S$ is a set of orthogonal idempotents in $R$ 
with $|S|=\aleph_{n-1}$ and $R$ is a $p$-boolean ring with $|R|=2^{\aleph_{n-1}}$ or $\beth_{n-1}$.\\
\indent If $S$ is an arbitrary set with $|S|=\aleph_{n-1}$, then $R= \text{Hom}_{\mathrm{Set}}(S, \bb{F}_p)$ is a $p$-boolean ring with $|R|=2^{\aleph_{n-1}}$ containing the set $\{\delta_{t=s}\}_{s \in S}\subset R$ of orthogonal idempotents.
\begin{claim}
There exists a $p$-boolean ring $R$ with $|R| = \beth_n$ and a subset $S \subset R$ consisting of orthogonal idempotents with $|S|=\beth_n$.
\end{claim}
\begin{proof}
We will prove the claim by induction on $n$. In the case $n=0$, we remark that the $p$-boolean ring $R_0=F_{\bb{F}_p}^p(R_0^{\oplus \bb{N}})=\bb{F}_p[X_1,X_2,...]/(X_1^p-X_1,X_2^p-X_2,...)$ admits a countably infinite subset of orthogonal idempotents given by a sequence $S_0=\{a_n\}_{n \in \bb{N}}$ of elements of $R$ defined by $a_n = (1-X^{p-1}_n) \prod_{0 \le i \le n-1} X^{p-1}_{i}$. The sequence $\{S_0, a: S_0 \to R_0\}$ solves the $n=0$ base-case.\\
\indent Suppose by induction we have a ring $R_{n-1}$ such that $|R_{n-1}| = \beth_{n-1}$ and that there is a subset $S_{n-1} \subset R_{n-1}$ consisting of orthogonal idempotents with $|S_{n-1}| = \beth_{n-1}$. Consider the ring:
\[R_n=\prod_{S_{n-1}}R_{n-1}.\]
We note $|R_n| = \beth_{n-1}^{\beth_{n-1}} = \beth_n$, and the set
\[S_n = \prod_{S_{n-1}} S_{n-1} \subset R_n\]
consists of orthogonal idempotents and $|S_n|=\beth_{n-1}^{\beth_{n-1}}=\beth_n$, as desired.
\end{proof}
\end{proof}
In Corollary~\ref{corollary:mainindiv} (ii), our cardinality bounds are optimal under the assumption of the Generalized Continuum Hypothesis (GCH). 
Indeed, if $f: A \to A'$ is a faithfully flat map of $p$-boolean, then we know that if $A$ is $\aleph_{n-1}$-counterable, 
then $f$ has descendability exponent at most $n$ by Theorem~\ref{thm:gj0}, so we would prefer to have upgraded our cardinality bounds from
$\text{min}(\beth_{n-1}, 2^{\aleph_{n-1}})$-countable to $\aleph_{n-1}$-countable.\\
\indent On the other hand, (i) is optimal; our ring $A$ is $\aleph_{n-1}$-Noetherian of Krull-dimension $n$, so again by Theorem~\ref{thm:gj0} 
any faithfully flat ring map out of it has descendability exponent at most $n$. However, the ring is quite pathological, since $A$ contains
a large collection of idempotents. In the next section, we will see how to construct more `geometric' examples of rings with faithfully flat maps of high descendability exponent.
\section{The case of an infinite polynomial ring}
The goal of this section is to construct a faithfully flat 
algebra over the infinite polynomial ring on an algebraically closed field $k[x_1,x_2,...]$ that is not descendable. The argument
will essentially boil down to a refinement of Theorem~\ref{theorem:indivisiblecup}, 
where we will need to construct a function $f:k \to k$ that is not polynomial on any infinite subset of $k$.\\
\indent We begin by recalling the following form of Lagrange's interpolation theorem.
\begin{theorem}[Lagrange interpolation]
Let $k$ be a field and $f \in k[x]$ a polynomial of degree $\le n$. Then there exists a polynomial $P_n \in \bb{Z}[x_1,...,x_n]$ such that for $z \in k^{\times n+1}$ we have:
\[\sum_{i=1}^{n+1} (-1)^iP_n(p^{n+1}_i(z))f(\text{pr}_i(z)) = 0.\]
Explicitly, $P = \prod_{ 1 \le i < j \le n} (x_j - x_i)$.
\end{theorem}
\begin{definition}
Let $k$ be a field. For $z \in k^{\times n+1}$, we may define a $k$-linear operator 
\[\nabla^z_i: \Hom(k^{\times m}, k) \to \Hom(k^{\times m-1},k)\]
as follows: For any function $f \in  \Hom(k^{\times m}, k)$, set:
\[\forall s \in k^{\times m-1}: \; \;  \nabla^{z}_i(f)(s) = \sum_{j=1}^{n+1} (-1)^j P_n(p^{n+1}_j(z)) f(t^{s}_i(\text{pr}_j(z))).\]
For any two $z \in k^{\times n+1}, z' \in k^{\times n'+1}$, one can verify that:
\[\nabla^z_i \nabla^{z'}_j = \nabla^{z'}_j \nabla^{z}_i\]
for $i \neq j$. Furthermore, if $f(x_1,...,x_m) = \prod_i f_i(x_i)$, and $z^i \in \Hom(k^{\times n_i}, k)$ for $i\in \{1,2,...,m\}$, 
\[(\nabla^{z^m}_m \circ \nabla^{z^{m-1}}_{m-1} \circ ... \circ \nabla^{z^1}_{1})(f) = \prod_i \nabla^{z^i} f_i.\]
\end{definition}
\begin{corollary}\label{corollary:interpol}
Let $f: \Omega \to k$ be function defined on a countably infinite subset $\Omega \subset k$. 
Then: $\nabla^{z} f = 0,\;  \forall z \in \Omega^{\times n+1}$ if and only if $f$ is a polynomial of degree $\le n$.
\end{corollary}
Corollary~\ref{corollary:interpol} provides a convenient criterion for checking if a function is polynomial on a given subset.
Using this, we can now construct functions on a field that are not polynomial on any infinite subset.
\begin{proposition}\label{proposition:notpoly}
For any field $k$, there exists an algebraically closed field extension $k \subset E(k)$, and a function $f: E(k) \to E(k)$ 
such that there is no countably infinite subset $S \subset E(k)$ such that $f$, when restricted to $S$, is a polynomial. 
If $|k| = \infty,$ then $|E(k)| = |k|$.
\end{proposition}
\begin{proof}
We define a sequence of field extensions $0=k_0\subset k=k_1 \subset k_2 \subset ...$ and functions $f_{n}:k_n \to k_{n+1}$ inductively. For the fields, set 
\[k_{n+1} = \overline{k_{n}(k_{n} \setminus k_{n-1})}\]
That is, $k_{n+1}$ is the algebraic closure of the field of fractions of the free polynomial ring $k_{n}[X_s]$ indexed by elements $s \in k_{n} \setminus k_{n-1}$. For the functions, we take:
\[\forall s \in k_{n} \setminus k_{n-1}: \; \; f_n(s) = X_s \in k_{n+1}\]
\[\forall s \in k_{n-1}: \; \; f_n(s) = f_{n-1}(s) \in k_{n} \subset k_{n+1}.\]
Set $E(k) = \underset{n}{\cup} k_n$, and $f=\colim_{n} f_n$, which is a map $E(k) \to E(k)$. \\
\indent We claim that $f$ is not polynomial on any infinite subset $S \subset E(k)$. Suppose that there exists a polynomial $p$ of degree $d$ such that $f|_{S} = p|_{S}$. There are two cases; there exists a $t \in \bb{N}$ such that $|S \cap k_t\setminus k_{t-1}| \ge d+1$, or for any $n \in \bb{N}$, $S \cap (k \setminus k_n) \neq 0$.\\
\indent In the first case, there exists an element $z \in (k_t \setminus k_{t-1})^{d+1}$ such that:
\[\nabla^z f = \sum_i (-1)^iP_d(p^{d+1}_i(z))f(\text{pr}_i(z))= 0.\]
As $P_d(p^{d+1}_i(z)) \in k(z_1,...,z_{n+1})$ for every $i$, this is a contradiction by construction.\\
\indent In the second case, take an $N \in \bb{N}$ such that the coefficients of $p$ are contained in $k_N$, and an $s \in S \cap (k \setminus k_N)$. Suppose that $s \in k_n\setminus k_{n-1}$. By supposition, $f(s) = p(s) \in k_n$, which is again a contradiction by construction.
\end{proof}
Later, we will need some results of infinite Ramsey theory, which originated from the paper \cite{ErdosRado}. Before stating the main theorem in a more convenient language, let us introduce some notation. 
\begin{construction}
Define the \textit{beth numbers} $\beth_n$ by the following formula:
\[\beth_0 = \aleph_0, \; \beth_n = 2^{\beth_{n-1}}.\]
\end{construction}
Note that $\beth_{n} \ge \aleph_{n}$, and equality holds if the generalized continuum hypothesis is assumed.
 For any set $S$, define $[S]^r=\{K \subset S: \; |K|=r\}$, for any cardinal $\alpha$, define $\alpha_{+}$ to be its 
 successor, and define $\beth_{\omega}=\bigcup_{\beta < \omega} \beth_{\beta}$ where $\omega$ is the first limit ordinal.
\begin{theorem}[Theorem 7.3 \cite{Higherinfinite}]\label{theorem:erdosrado}
For any set $S$ of size $\beth_{r+}$, and any set map:
\[f: [S]^{r+1} \to \bb{N}\]
There exists a set $S' \subset S$ with $|S'|=\aleph_{1}$ and an integer $n \in \bb{N}$ such that $[S']^{r+1} \subset f^{-1}(n)$.
\end{theorem}
\begin{remark}\label{remark:boxuseful}
Theorem~\ref{theorem:erdosrado} also implies the following weaker statement: For a set $S$ of cardinality $\beth_{r+}$, and a set theoretic map
\[f: S^{\times r+1} \to \bb{N}\]
such that $f(s) = f(\sigma(s))$ for every $s \in S^{\times r+1}$ and $\sigma \in \mathrm{S}_{r+1}$, the symmetric group on $r+1$ elements, there exists an $n \in \bb{N}$ and countably infinite subsets $S_1, ..., S_{r+1}$ of $S$ such that $ \prod_{i=1}^{r+1} S_i \subset f^{-1}(n)$.
\end{remark}
Armed with these tools, we may now state and prove our refinement of Theorem~\ref{theorem:indivisiblecup}.
\begin{theorem}\label{theorem:countermod}
For $m \in \bb{N} \cup \{\infty\}$, let $F$ be a field of cardinality $\beth_{m-1+}$, and $k=E(F)$ (see Proposition~\ref{proposition:notpoly}). 
Let $R_m=k[x_1,...,x_m]$, and for any $n < m+1$, let $R_n=k[x_1,...,x_n]$ and $\{\{S_i,\Psi_{i}: k \to R_n\}\}$ be the $n$-indivisible sequence 
of Example~\ref{example:indiv} (i) (where $S_i=k$). Let $i:R_n \to R_m$ be the canonical inclusion, and $\{\{S_i,i \circ \Psi_{i}: k \to R_m\}\}$ be 
the induced $n$-indivisible sequence in $R_m$ (see Proposition~\ref{proposition:indivequiv}).
Then there exists an $k[x]$-linear map $h: \oplus_{k} k[x] \to k[x]$ such that:
\[\bigotimes_{i\in [n], R_m} (h \otimes_{k[x], t_i} R_m \circ \cl{\Psi_{R_m,S_i}}) \neq 0 \in \Ext^n_{R_m}(\otimes_{i\in [n], R_m} \coker(\Psi_{R_m,S_i}), R_m),\]
where $t_i: k[x] \to R_m$ is the $k$-algebra map sending $x$ to $x_i$.
\end{theorem}
\begin{proof}
The map $R_n \to R_m$ admits an $R_n$-linear left-inverse, and since 
\[\cl{\Psi_{R_m,S_i}} = \cl{\Psi_{R_n,S_i}} \otimes_{R_n} R_n, \]
it suffices to show the claim for $m=n$. By Lemma~\ref{lemma:extcrit}, we need to show that there is no $R_n$-linear extension $r$ fitting into the diagram:
\[\begin{tikzcd} R_n & \\
\bigoplus_{k^{\times n}} R_n \arrow[r, "\oplus_i \Psi_{R_n,S_i}"] \arrow[u, "h^{\otimes_k n}"] & \bigoplus_{i} \bigoplus_{k^{\times n-1}} R_n \oplus \bigoplus_{k^{\times n}} R_n \arrow[ul, swap, "r", dashed]\end{tikzcd}\]
We shall assume such an $r$ exists and derive a contradiction.\\
\indent Let $H: k \to k$ be a function that is not polynomial on any countably infinite subset of $k$ (see Proposition~\ref{proposition:notpoly}). Set:
\[h(1_{s}) = H(s) \in k[x],\]
for every $s \in k$. Further, for every $s' \in k^{\times n-1}$ and $s \in k^{\times n}$, let
\[r(1_{i,s'}) = f_i(s'), r(1_{i,s}) = g_i(s).\]
Under the identification $R_n = k[x_1,...,x_n]$, $f_i(s'), g_i(s)$ are polynomials subject to the equation:
\[\sum^n_{i=1} f_i(p^n_i(t))(x_1,...,x_n) + \sum^n_{i=1} (x_i - \text{pr}_i(t)) g_i(t)(x_1,...,x_n) = \prod^n_{i=1} h(1_{\text{pr}_i(t)})(x_i).\]
for every $t \in k^{\times n}$. In particular, we have:
\[\sum^n_{i=1} f_i(p^n_i(t))(\text{pr}_1(t),...,\text{pr}_n(t)) = \prod^n_{i=1} h(1_{\text{pr}_i(t)})(\text{pr}_i(t)). \; \; (*)\]
\begin{claim}
There exists a countably infinite subsets $S_1, ..., S_n$ of $k$, such that for every $\omega \in \Omega =\prod^n_i S_i$, the polynomials $f_i(p^n_i(\omega))$ have (total) degree $\le m$ for some $m \in \bb{N}$.
\end{claim}
\begin{proof}
Consider the function:
\[c: k^{\times n} \to \bb{N}\]
which takes a point $t \in k^{\times n}$ to $\underset{i\in\{1,2,...,n\}}{\text{max}}(\text{deg}(f_i(p^n_i(t)))$. By Remark~\ref{remark:boxuseful}, there exists a $m \in \bb{N}$ and countably infinite subsets $S_1, ..., S_n$ such that $ \prod_{i=1}^{n} S_i \subset f^{-1}(m)$. Set $\Omega = \prod_{i=1}^n S_i$ to get the claim.
\end{proof}
For each $i \in \{1,2,...,n\}$, let us take points $z^i \in S_i^{\times m+1}$. Note that for all $i \in \{1,2,...,n\}$, we have functions $\tilde{f}_i, \tilde{h} \in \Hom_{\mathrm{Set}}(k^{\times n},k)$ defined as follows:
\[\forall t \in k^{\times n}: \; \; \tilde{f}_i(t) = f_i(p^n_i(t))(\text{pr}_1(t),...,\text{pr}_n(t))\]
\[\forall t \in k^{\times n}: \; \; \tilde{h}(t) = \prod_{i=1}^n H(\text{pr}_i(t)).\]
We see that $(*)$ may be re-written as an equality
\[\sum^n_{i=1} \tilde{f}_i = \tilde{h},\]
of functions in $\Hom_{\mathrm{Set}}(k^{\times n}, k)$. For any $s \in \prod_{j \in \{1,2,...,n\} \setminus \{i\}} S_j$, we see
\[\tilde{f}_i \circ t^{s}_i: k \to k\]
is a polynomial of degree $\le m$. Therefore:
\[\nabla^{z^i}_i \tilde{f}_i = 0,\]
for all $i \in \{1,2,...,n\}$. Therefore:
\[ \prod_i^n \nabla^{z^i} H = (\nabla^{z^n}_n \circ \nabla^{z^{n-1}}_{n-1} \circ ... \circ \nabla^{z^1}_1)( \tilde{h}) = (\nabla^{z^n}_n \circ \nabla^{z^{n-1}}_{n-1} \circ ... \circ \nabla^{z^1}_1)(\sum_{i=1}^n \tilde{f}_i)=0.\]
This implies that there exists an $i \in \{1,2,...,n\}$ such that $\nabla^{z^i} H = 0$ for any such $z^i \in S_i^{\times m+1} \subset k^{m+1}$. But this would imply that $H$ is a polynomial of degree $m$ on the infinite subset $S_i \subset k$ (due to Corollary~\ref{corollary:interpol}), which is a contradiction.
\end{proof}
\begin{theorem}\label{theorem:maincounter}
$F$ a field with $|F| = \beth_{n-1+}$, and $k=E(F)$. 
For any $n \in \bb{N} \cup \{\infty\}$, there exists a faithfully flat ring map $k[x_1,...,x_n] \to A_{n}$ 
whose descendability exponent is $n$. In particular, setting $n=\infty$, we obtain a faithfully flat ring map $k[x_1,x_2,...] \to A_{\infty}$ 
that is not descendable.
\end{theorem}
\begin{proof}
For $n \in \bb{N} \cup \{\infty\}$, let $R=k[x_1,...,x_n]$, let $\{\{S_i, \Psi_{i}: k \to A\}\}$ be the $n$-indivisible sequence of Example~\ref{example:indiv} (i) (where $S_i=k$). 
Let $h: \oplus_{k} k[x] \to k[x]$ be the map constructed in Theorem~\ref{theorem:countermod}, and consider the commutative diagram:
\[
\begin{tikzcd}[column sep=6em, row sep=4em]
R \arrow[r] & A_n \\
\mathrm{Sym}_{R}(\oplus_{S_i} R)
  \arrow[u, "g_i"', swap]
  \arrow[r, "\mathrm{Sym}_{R}(\Psi_{R,S_i})"]
& \mathrm{Sym}_{R}(R \oplus \bigoplus_{S_i} R)
  \arrow[u] \\
\oplus_{S_i} R
  \arrow[u]
  \arrow[uu, bend left=50, looseness=1.4, "h^{\otimes_{k} n}"', swap]
  \arrow[r, "\Psi_{R,S_i}"]
& R \oplus \bigoplus_{S_i} R
  \arrow[u]
\end{tikzcd}
\]
where $g_i$ is the $R$-algebra map induced by $h^{\otimes_{k} n}$, and 
\[A_n = R \otimes_{g_i,\mathrm{Sym}_{R}(\oplus_{S_i} R),\mathrm{Sym}_{R}(\Psi_{R,S_i})} \mathrm{Sym}_{R}(R \oplus \bigoplus_{S_i} R).\]
By Proposition~\ref{proposition:symformula}, the map $\mathrm{Sym}_{R}(\Psi_{R,S_i})$ is a faithfully flat map of commutative rings, 
so $R \to A_n$ is also faithfully flat. 
By Theorem~\ref{theorem:countermod}, the conditions of Proposition~\ref{proposition:modalg} are satisfied, so the 
descendability exponent of $R \to A_n$ is at least $n$. Setting $n=\infty$, we obtain a faithfully flat ring map $k[x_1,x_2,...] \to A_\infty$ that is not descendable.
\end{proof}



